\scanpage{4}
% archive foliation: 3
% The typescript is upside down relative to the scan; transcribed in its reading orientation.
\begin{center}-- 39 --\end{center}
\hfill n$^{\circ}$329

point $a$, $f_{\unsure{x|\varkappa}}$ est continue au point $a_{\unsure{x|\varkappa}}$.

COROLLAIRE 2. -- \uline{Soient} $X$, $Y$ \uline{deux espaces topologiques. Pour qu'une application} $f:X\to Y$ \uline{soit continue, il faut et il suffit que l'application} $g:x\to(x,f(x))$ \uline{soit un homéomorphisme de} $X$ \uline{sur le graphe} $G$ \uline{de} $f$ (\uline{considéré comme sous-espace de l'espace produit} $X\times Y$).

Comme $f=\operatorname{pr}_2\circ g$, la condition est suffisante. Elle est nécessaire, car si $f$ est continue, $g$ est bijective et continue (cor. 1) et son application réciproque est la restriction à $G$ de $\operatorname{pr}_1$, qui est continue (cf. \uline{Ens.}, Chap. IV, § 2, n$^{\circ}$4, critère CST 17).

\uline{PROPOSITION 2.} (associativité des produits topologiques). -- \uline{Soient} $(X_i)_{i\in I}$ \uline{une famille d'espaces topologiques,} $(J_{\unsure{x|\varkappa}})_{\unsure{x|\varkappa}\in K}$ \uline{une partition de l'ensemble} $I$, \uline{et, pour tout} $\unsure{x|\varkappa}\in K$, \uline{soit} $X'_{\unsure{x|\varkappa}}=\prod_{i\in J_{\unsure{x|\varkappa}}}X_i$ \uline{le produit des espaces} $X_i$ \uline{pour} $i\in J_{\unsure{x|\varkappa}}$. \uline{L'application canonique} (\uline{Ens.} R, § 4, n$^{\circ}$ 11) \uline{de l'espace produit} $\prod_{i\in I}X_i$ \uline{sur l'espace produit} $\prod_{\unsure{x|\varkappa}\in K}X'_{\unsure{x|\varkappa}}$ \uline{est un homéomorphisme.}

C'est un cas particulier de la transitivité des topologies initiales (§ 2, n$^{\circ}$ 3, prop. 5, cf. \uline{Ens.}, chap. IV, § 2, n$^{\circ}$4, critère CST 13).

On \uline{identifie} le plus souvent les espaces produits $\prod_{i\in I}X_i$ et $\prod_{\unsure{x|\varkappa}\in K}X'_{\unsure{x|\varkappa}}$ au moyen de l'application canonique.

COROLLAIRE. -- \uline{Soit} $\sigma$ \uline{une permutation de l'ensemble} $I$. \uline{L'application} $(x_i)\to(x_{\sigma(i)})$ \uline{est un homéomorphisme de} $\prod_I X_i$ \uline{sur} $\prod_{i\in I}X_{\sigma(i)}$

Il suffit, dans la prop.2, de prendre $K=I$ et $J_i=\{\sigma(i)\}$ pour tout $i\in I$.

PROPOSITION 3. -- \uline{Soient} $X$ \uline{un ensemble,} $(Y_i)_{i\in I}$ \uline{une famille d'espaces topologiques, et pour chaque} $i\in I$, \uline{soit} $f_i$ \uline{une application de} $X$ \uline{dans} $Y_i$ \uline{Soit} $f$ \uline{l'application} $x\to(f_i(x))$ de $X$ \uline{dans} $Y=\prod_{i\in I}Y_i$, \uline{et soit} $\mathcal{T}$ \uline{la topologie la moins fine sur} $X$ \uline{rendant continues les} $f_i$. \uline{Alors} $\mathcal{T}$ \uline{est}
